\chapter{Error Handling and Deferred Execution}
\label{ch:error-defer}

Real programs deal with things that can go wrong: a withdrawal larger than the
balance, a number that will not parse, a file that is not there. Equally, real
programs acquire resources open files, allocated memory that must be
released no matter how the work turns out. This chapter covers Salam's
straightforward approach to both: reporting failure through return values, and
guaranteeing cleanup with \code{defer}.

\section{Salam's philosophy: errors are values}

Salam has no exceptions, no \code{try}/\code{catch}, no invisible control flow
that whisks you off to some distant handler. Instead, a function that can fail
\emph{says so in what it returns}, and the caller checks. This keeps the flow of
your program on the page: every place that can fail is visible, and every failure
is handled where it happens. There are a few common patterns, from simplest to
most expressive.

\section{Pattern 1: a boolean success flag}

The simplest signal is a \code{bool}: \code{true} for success, \code{false} for
failure. You saw this in Chapter~9 with the bank account, whose \code{withdraw}
returned whether it succeeded:

\begin{salam}
struct Account:
    balance : int = 0
    func withdraw(amount : int) : bool:
        if amount > this.balance:
            ret false              // failure: not enough funds
        end
        this.balance = this.balance - amount
        ret true                   // success
    end
end

func main:
    mut acc : Account = Account { balance = 100 }
    if acc.withdraw(30):
        println "withdrew 30; balance is", acc.balance
    else:
        println "withdrawal failed"
    end
    if acc.withdraw(1000):
        println "withdrew 1000"
    else:
        println "withdrawal failed: insufficient funds"
    end
end
\end{salam}

\begin{progout}
withdrew 30; balance is 70
withdrawal failed: insufficient funds
\end{progout}

The caller cannot ignore the outcome by accident: it is right there in the
\code{if}. This pattern is perfect when the only thing you need to know is
``did it work?''

\section{Pattern 2: a sentinel return value}

When a function returns a value but might have nothing meaningful to give back, a
\textbf{sentinel} a special value that means ``no result'' is a common
choice. A search that returns the index of an item can return \code{-1} to mean
``not found,'' since a real index is never negative:

\begin{salam}
func index_of(values : int[5], target : int) : int:
    for mut i = 0, i < 5, i = i + 1:
        if values[i] == target:
            ret i
        end
    end
    ret -1                         // sentinel: not found
end

func main:
    nums : int[5] = [10, 20, 30, 40, 50]
    found : auto = index_of(nums, 30)
    if found >= 0:
        println "found at index", found
    else:
        println "not found"
    end
    println "search for 99:", index_of(nums, 99)
end
\end{salam}

\begin{progout}
found at index 2
search for 99: -1
\end{progout}

\begin{warn}
A sentinel only works when there is a value the real results can never take. For
an index, \code{-1} is safe. But ``zero means failure'' is dangerous if zero is a
legitimate result, and ``empty string means error'' fails if the empty string is
valid data. When no value can be safely reserved, prefer the \code{Option}
approach below.
\end{warn}

\section{Pattern 3: \texttt{Option} a value that may be absent}

The standard library's \code{Option} type makes ``maybe a value, maybe nothing''
explicit and safe, without stealing a value from the result's range. An
\code{Option<T>} either holds a \code{T} (it is \emph{some} value) or holds
nothing (it is \emph{none}). Import it with \code{import "option"}.

\begin{salam}
import "option"

func safe_divide(a : int, b : int) : option.Option<int>:
    if b == 0:
        ret Option { value = 0, present = false }    // none
    end
    ret option.Some(a / b)                            // some result
end

func main:
    r1 : auto = safe_divide(10, 2)
    r2 : auto = safe_divide(10, 0)
    println "10/2 is some?", option.IsSome(r1)
    println "10/0 is some?", option.IsSome(r2)
    println "10/2 =", option.UnwrapOr(r1, -1)
    println "10/0 =", option.UnwrapOr(r2, -1)    // falls back to -1
end
\end{salam}

\begin{progout}
10/2 is some? true
10/0 is some? false
10/2 = 0
10/0 = -1
\end{progout}

The \code{option} module gives you the tools to work with these values safely:

\begin{center}
\begin{tabular}{@{}ll@{}}
\toprule
\textbf{Function} & \textbf{What it does} \\
\midrule
\code{option.Some(x)}        & wrap \code{x} as a present value \\
\code{option.IsSome(o)}      & \code{true} if \code{o} holds a value \\
\code{option.IsNone(o)}      & \code{true} if \code{o} holds nothing \\
\code{option.UnwrapOr(o, d)} & the held value, or \code{d} if none \\
\bottomrule
\end{tabular}
\end{center}

\code{UnwrapOr} is especially handy: it gets the value out, supplying a fallback
for the ``none'' case in one step, so you never accidentally use a value that is
not there. To build a ``none,'' you construct the \code{Option} with
\code{present = false}, as the example shows.

\begin{note}
\code{Option} is the type to reach for whenever ``no answer'' is a legitimate
outcome: a lookup that may miss, a parse that may fail, the first element of a
possibly-empty collection. It forces the caller to acknowledge the empty case,
turning a whole class of ``forgot to check'' bugs into something the type system
helps you remember.
\end{note}

\section{Deferred execution with \texttt{defer}}

Now to the second half of the chapter: making sure cleanup happens. The
\code{defer} keyword schedules a statement to run \emph{when the current function
exits} no matter how it exits. You write the cleanup right next to the thing it
cleans up, and Salam guarantees it runs later:

\begin{salam}
func main:
    defer println "cleanup runs last"
    println "doing work"
    println "more work"
end
\end{salam}

\begin{progout}
doing work
more work
cleanup runs last
\end{progout}

The \code{defer}'d line was \emph{written} first but \emph{ran} last, at the
moment \code{main} finished. The point is that the cleanup is registered the
instant you set up whatever needs cleaning, so you can never forget it later.

\subsection{Defers run in reverse order}

If you register several defers, they run in \textbf{last-in, first-out} order ---
the most recently deferred runs first. This mirrors how resources nest: the last
thing you opened is the first thing you should close.

\begin{salam}
func main:
    defer println "cleanup 1"
    defer println "cleanup 2"
    println "start"
    println "end"
end
\end{salam}

\begin{progout}
start
end
cleanup 2
cleanup 1
\end{progout}

\subsection{Defer runs even on an early return}

The real power of \code{defer} is that it runs however the function ends ---
including an early \code{ret} from somewhere in the middle. This is what makes it
reliable: the cleanup cannot be skipped by a code path you forgot about.

\begin{salam}
func greet(name : str):
    defer println "done greeting"
    println "hello"
    if name == "":
        println "no name given"
        ret                        // early return -- defer still runs
    end
    println name
end

func main:
    greet("Ali")
    println "---"
    greet("")
end
\end{salam}

\begin{progout}
hello
Ali
done greeting
---
hello
no name given
done greeting
\end{progout}

Both calls print \code{"done greeting"} the normal path through \code{Ali} and
the early-return path through the empty name. You write the cleanup once, and it
fires on every exit.

\section{Putting it together: acquire, defer, use}

The idiomatic pattern is: acquire a resource, immediately \code{defer} its
release, then use it freely without worrying about the exits. Whether the function
returns normally, returns early on an error, or falls through, the resource is
released exactly once. We will see this with real files and memory in Chapters~17
and~18; the shape is always the same:

\begin{salam}
func process():
    // acquire something here
    defer println "release the resource"

    println "use the resource"
    if true:
        println "early exit on some condition"
        ret               // the defer above still runs
    end
    println "this line is skipped"
end

func main:
    process()
end
\end{salam}

\begin{progout}
use the resource
early exit on some condition
release the resource
\end{progout}

\begin{tip}
Pair every acquisition with a \code{defer} on the very next line. Reading top to
bottom, ``open file; defer close file'' is impossible to get wrong the close is
guaranteed and sits right where you can see it. This single habit eliminates the
most common kind of resource leak.
\end{tip}

\section{Chapter summary}

\begin{itemize}[leftmargin=1.4em]
  \item Salam reports errors through \emph{return values}, not exceptions every
        failure is visible and handled at the call site.
  \item Return a \code{bool} when the caller only needs ``did it work?''.
  \item Return a \textbf{sentinel} (like \code{-1}) when a value's range has room
        for a ``no result'' marker but only then.
  \item Use \code{option.Option<T>} (\code{Some}, \code{IsSome}, \code{IsNone},
        \code{UnwrapOr}) when ``no value'' must be represented safely.
  \item \code{defer} schedules a statement to run when the function exits, in
        last-in-first-out order, \emph{including} on early returns.
  \item The reliable pattern is: acquire, immediately \code{defer} the release,
        then use.
\end{itemize}

\section{Exercises}

\exercise Write \code{safe\_sqrt(x f64) option.Option<f64>} that returns
\code{none} for negative input and \code{Some} of the square root otherwise (use
the \code{math} module from the next chapter, or \code{x ** 0.5}).

\exercise Write a function \code{find\_first\_negative(a int[5]) int} that returns
the index of the first negative element or \code{-1} if there is none.

\exercise Add two \code{defer} statements and a normal statement to a function,
and predict the output order before running it.

\exercise Write a function with an early \code{ret} guarded by a condition and a
\code{defer} that prints ``finished''. Verify ``finished'' prints on both paths.

\exercise Rewrite the \code{safe\_divide} example so the caller prints a friendly
message in the ``none'' case instead of using \code{UnwrapOr}.
